% Eigenständiger Prosabeweis; die nummerierten Aussagen bleiben in Band 40.
\section*{Kann man die ursprünglichen Elemente noch erkennen?}
\hypertarget{reconstruction.reading.question}{}

Eine Gruppe besteht aus Elementen, die man miteinander multiplizieren
kann. Man kann aber auch \emph{Mengen} von Gruppenelementen
multiplizieren: Man bildet alle Produkte, die durch eine Wahl aus
der ersten und eine Wahl aus der zweiten Menge entstehen.
Für zwei nichtleere Teilmengen \(X,Y\) einer Gruppe \(G\) setzen wir
\[
  XY=\{xy\mid x\in X,\ y\in Y\}.
\]
So enthält \(\{a,b\}\{c,d\}\) die Produkte \(ac,ad,bc,bd\).
Verschiedene Produkte können denselben Wert haben; das Ergebnis
muss daher keine vier Elemente besitzen.

Nun stellen wir uns vor, dass nur noch diese größere Rechenwelt
bekannt ist: ihre Eingaben sind sämtliche nichtleeren Teilmengen,
ihre Operation ist das Mengenprodukt. Die Beschriftungen, an denen
man Einermengen oder die Anzahl der enthaltenen Elemente erkennt,
sind verloren gegangen. Lässt sich die ursprüngliche Gruppe
trotzdem zurückgewinnen?

Diese Frage ist anspruchsvoller, als sie zunächst klingt. Zwar
gilt für Einermengen immer
\[
  \{g\}\{h\}=\{gh\}.
\]
Könnten wir die Einermengen erkennen, hätten wir die ursprüngliche
Multiplikation schon vor uns. Aber genau ihre Erkennbarkeit müssen
wir begründen. Ein Isomorphismus der größeren Rechenwelten erhält
die Multiplikation; dass er auch Mengengrößen oder Inklusionen
erhält, ist keine Voraussetzung.

Wir werden zwei Aussagen beweisen. Zunächst zeigen wir: Zwei
Gruppen mit isomorphen Potenzhalbgruppen sind selbst isomorph.
Danach gehen wir einen Schritt weiter: Ist nur eine Ausgangsstruktur
als Gruppe bekannt und die andere zunächst bloß eine Halbgruppe,
so muss auch die zweite eine Gruppe sein. Der erste Beweis beruht
auf Einermengen, der zweite zusätzlich auf der gesamten Menge der
invertierbaren Elemente.

\medskip\noindent\textbf{Was bedeutet hier zurückgewinnen?}
Gesucht sind keine ursprünglichen Namen der Elemente. Gesucht ist
eine Gruppe mit derselben Multiplikation bis auf Umbenennung.
Zwei so übereinstimmende Gruppen heißen \emph{isomorph}. Wenn jede
Gruppe mit derselben Potenzhalbgruppenstruktur zu \(G\) isomorph
sein muss, ist \(G\) in genau diesem Sinn eindeutig bestimmt.
Noch konkreter werden wir innerhalb der großen Rechenwelt eine
kleinere Gruppe auswählen: ihre invertierbaren Elemente. Dass
diese gerade den ursprünglichen Einermengen entsprechen, erklärt,
warum unsere beiden Sätze die Eingangsfrage beantworten.

\par\Needspace{14\baselineskip}
\section*{Die Zweiergruppe als vollständiges Beispiel}
\hypertarget{reconstruction.reading.example}{}

Wir betrachten die Gruppe \(G=\{e,s\}\) mit neutralem Element
\(e\) und \(s^2=e\). Ihre drei nichtleeren Teilmengen sind
\(\{e\}\), \(\{s\}\) und \(G\). Für ihre Mengenprodukte ergibt
sich die vollständige Tafel
\[
  \renewcommand{\arraystretch}{1.35}
  \begin{array}{c|ccc}
    \cdot&\{e\}&\{s\}&G\\\hline
    \{e\}&\{e\}&\{s\}&G\\
    \{s\}&\{s\}&\{e\}&G\\
    G&G&G&G
  \end{array}
\]
Die Zeile \(\{s\}\), Spalte \(G\), enthält beispielsweise
\(\{se,ss\}=\{s,e\}=G\). Beim Produkt \(GG\) entstehen
\(ee,es,se,ss\), also ebenfalls genau \(e\) und \(s\).

Die Einermenge \(\{e\}\) ist das neutrale Element dieser Tafel:
Multiplikation mit ihr verändert keine Eingabe. Die Einermenge
\(\{s\}\) besitzt ein zweiseitiges Inverses, nämlich sich selbst.
Auch \(\{e\}\) ist ihr eigenes Inverses. Für \(G\) gibt es dagegen
kein Inverses: Jeder Eintrag seiner Zeile und seiner Spalte ist
\(G\), niemals \(\{e\}\).

Damit lässt sich in diesem Beispiel die ursprüngliche Gruppe
an einer reinen Recheneigenschaft erkennen. Die beiden
\emph{invertierbaren} Eingaben bilden genau die Einermengen.
Hießen die drei Eingaben nur \(p,q,r\), könnten wir immer noch
das neutrale Element bestimmen und anschließend diejenigen
Eingaben auswählen, die sich zu ihm zurückmultiplizieren lassen.
Wir müssten nicht wissen, welche Mengen hinter den Namen stehen.

Das Vorgehen lässt sich an der unbeschrifteten Tafel Schritt für
Schritt ausführen. Zuerst suchen wir die Eingabe \(p\), deren Zeile
und Spalte jede Eingabe unverändert lassen. Dann wählen wir alle
Eingaben \(q\), für die es eine Eingabe \(r\) mit \(qr=rq=p\) gibt.
Schließlich behalten wir nur die Zeilen und Spalten dieser
invertierbaren Eingaben. Im Beispiel bleibt die Tafel
\[
  \begin{array}{c|cc}
    \cdot&p&q\\\hline
    p&p&q\\
    q&q&p
  \end{array}
\]
übrig. Die Zuordnung \(e\mapsto p\), \(s\mapsto q\) erhält
alle vier Produkte: Wir haben die Zweiergruppe wiedergefunden.
Bei unendlichen Gruppen ist dies eine mathematische Beschreibung,
kein Versprechen eines endlichen Suchverfahrens. Warum die
ausgewählte Schicht stets genau die richtige ist, beweisen wir nun.

Der entscheidende Punkt ist, dass derselbe Gedanke für jede
Gruppe funktioniert, auch für unendliche und nichtkommutative
Gruppen. Um später auch die zweite Beweisstufe zu verstehen,
formulieren wir ihn gleich etwas allgemeiner.

\par\Needspace{13\baselineskip}
\section*{Invertierbare Mengen sind Einermengen}
\hypertarget{reconstruction.reading.units}{}

Eine \emph{Halbgruppe} ist eine Menge mit einer
assoziativen Multiplikation. Ein \emph{Monoid} ist eine Halbgruppe
mit neutralem Element \(e\). Ein Element \(u\) eines Monoids heißt
\emph{Einheit}, wenn es ein \(v\) mit
\[
  uv=vu=e
\]
gibt. Man nennt \(v\) dann das Inverse von \(u\) und schreibt
\(u^{-1}\). Das Inverse ist eindeutig: Sind \(v\) und \(w\)
beide invers zu \(u\), so ist
\(v=v(uw)=(vu)w=w\). Eine Gruppe ist ein Monoid, in dem
jedes Element eine Einheit ist.

Für eine Halbgruppe \(A\) bezeichnen wir die Menge ihrer
nichtleeren Teilmengen durch
\[
  \PowNE A=\{X\subseteq A\mid X\neq\varnothing\}.
\]
Das Mengenprodukt macht daraus wieder eine Halbgruppe. Das
Produkt zweier nichtleerer Teilmengen ist nichtleer und liegt
weiterhin in \(A\). Die Assoziativität folgt aus derjenigen von
\(A\), denn
\[
  \begin{aligned}
    (XY)Z&=\{(xy)z\mid x\in X,y\in Y,z\in Z\}\\
         &=\{x(yz)\mid x\in X,y\in Y,z\in Z\}=X(YZ).
  \end{aligned}
\]
Hinter dem mittleren Gleichheitszeichen steht ein allgemeines
Mengenprinzip: Stimmen zwei Ausdrücke bei jeder zulässigen Wahl
ihrer Variablen überein, erzeugen sie dieselbe Menge von Werten.
In Band~3 ist dies bereits als
\FormulaRefAuto{B27TermImagePointwiseEquality}[thm] mit einer
Beweistabelle festgehalten. Für unsere drei Variablen kann man
mit \(J=(X\times Y)\times Z\) indizieren: Zu jedem
\(((x,y),z)\in J\) gehören die beiden Werte \((xy)z\) und
\(x(yz)\), die wegen der Assoziativität gleich sind.
Das zusätzliche Dreifachzeugenschema
\FormulaRefAuto{TripleTermWitnessPointwiseEquality}[thm] und
sein Mengenschluss
\FormulaRefAuto{TripleTermRepresentationEquality}[thm]
formulieren diesen Schritt direkt für \(x,y,z\).
Damit braucht die zugehörige Tabelle in Band~28 die
Mengenargumentation nicht jeweils neu auszuführen.

Wir nennen sie die \emph{große Potenzhalbgruppe}: Zugelassen
sind alle nichtleeren Teilmengen, auch unendliche. Ist \(A\)
ein Monoid, so ist \(\{e\}\) neutral, weil
\(\{e\}X=X=X\{e\}\).

\medskip\noindent\textbf{Der Schlüssel.}
Ist \(A\) ein Monoid, so sind die Einheiten von \(\PowNE A\)
genau die Einermengen
\(\{u\}\), deren einziges Element \(u\) eine Einheit von \(A\)
ist. Dafür brauchen wir keine Kürzbarkeit des gesamten Monoids.

Die eine Richtung ist sofort sichtbar. Besitzt \(u\) das
Inverse \(v\), so gilt
\[
  \{u\}\{v\}=\{e\}=\{v\}\{u\}.
\]
Damit ist \(\{u\}\) invertierbar. Für die andere Richtung sei
\(X\) eine invertierbare nichtleere Teilmenge. Dann gibt es eine
nichtleere Teilmenge \(Y\) mit
\begin{equation}
  XY=\{e\}=YX.
  \tag{E}\label{reconstruction-reading-inverse}
\end{equation}
Wählen wir \(u\in X\) und \(v\in Y\), so folgt bereits
\(uv=vu=e\). Wir müssen noch zeigen, dass \(X\) und \(Y\)
keine weiteren Elemente enthalten können.

Sei \(x\in X\). Wegen \(XY=\{e\}\) ist auch \(xv=e\).
Mit Assoziativität und den Neutralitätsgesetzen folgt
\[
  x=xe=x(vu)=(xv)u=eu=u.
\]
Somit ist \(X=\{u\}\). Für \(y\in Y\) erhalten wir ebenso
aus \(uy=e\) die Gleichung
\[
  y=ey=(vu)y=v(uy)=ve=v.
\]
Also ist \(Y=\{v\}\), und \(u\) ist eine Einheit von \(A\).
Das beweist beide Richtungen. Die allgemeine Aussage steht
in Band~38 unter \FormulaRefAuto{PowerMonoidUnitCharacterization}[thm].

Für eine Gruppe \(G\) ist jedes \(g\in G\) invertierbar.
Deshalb erhalten wir genau
\begin{equation}
  \text{Einheiten von }\PowNE G
    =\{\{g\}\mid g\in G\}.
  \tag{S}\label{reconstruction-reading-singletons}
\end{equation}
Das ist \FormulaRefAuto{PowerGroupUnitsAreSingletons}[thm]
in Band~40. Die rechte Seite beschreibt die Einermengen mithilfe
der ursprünglichen Elemente. Die linke Seite erkennt dieselben
Objekte allein an der Multiplikation der Potenzhalbgruppe.
Genau diese zweite Beschreibung übersteht eine Umbenennung
der größeren Rechenwelt.

Damit ist auch das allgemeine Vorgehen begründet: Das neutrale
Element und die Einheiten werden allein durch Produktgleichungen
erkannt. Auf den Einheiten behalten wir die vorhandene
Multiplikation bei. Die Abbildung \(g\mapsto\{g\}\) ist eine
Bijektion auf diese Schicht und erhält wegen
\(\{g\}\{h\}=\{gh\}\) die Multiplikation. Die so ausgewählte
Gruppe ist deshalb eine isomorphe Kopie der Ausgangsgruppe.

\par\Needspace{12\baselineskip}
\section*{Aus dem großen Isomorphismus wird ein kleiner}
\hypertarget{reconstruction.reading.groups}{}

Seien zunächst \(G\) und \(H\) beide Gruppen und
\[
  F\colon\PowNE G\longrightarrow\PowNE H
\]
ein Isomorphismus ihrer Potenzhalbgruppen. Das bedeutet:
\(F\) ist bijektiv und für alle nichtleeren \(X,Y\subseteq G\)
gilt
\[
  F(XY)=F(X)F(Y).
\]
Eine Bijektion ordnet jedem Zielobjekt genau ein Urbild zu.
Insbesondere besitzt \(F\) eine Umkehrabbildung \(F^{-1}\).
Auch diese erhält Produkte: Schreiben wir \(Z=F(X)\) und
\(W=F(Y)\), so ist
\(F^{-1}(ZW)=F^{-1}(F(XY))=XY\).
Da \(X=F^{-1}(Z)\) und \(Y=F^{-1}(W)\), ist dies genau
\[
  F^{-1}(ZW)=F^{-1}(Z)F^{-1}(W).
\]
Die allgemeine Beweistabelle steht bereits in Band~28 unter
\FormulaRefAuto{SemigroupIsoInverse}[thm]. Bei Monoiden wird
zusätzlich das neutrale Element erhalten; bei Gruppen damit
auch jedes inverse Paar. Die ausdrücklichen Umkehrsätze sind
\FormulaRefAuto{MonoidIsoInverse}[thm] in Band~38 und
\FormulaRefAuto{GroupIsoInverse}[thm] in Band~40.

Dabei ist die Ebene der Aussage entscheidend: Aus einem
Isomorphismus \(G\to S\) mit \(G\) einer Gruppe folgt direkt,
dass \(S\) eine Gruppe ist
(\FormulaRefAuto{SemigroupIsoGroupTransport}[thm]). Hier
haben wir jedoch erst einen Isomorphismus
\(\PowNE G\to\PowNE S\). Ein Grundisomorphismus \(G\to S\)
muss daraus zunächst gewonnen werden. Genau das leistet
die Rekonstruktion; der allgemeine Umkehrsatz nimmt sie
nicht vorweg.

Zunächst muss \(F\) das neutrale Element auf das neutrale
Element abbilden. Für jedes \(Z\in\PowNE H\) wählen wir
sein Urbild \(X\). Dann ist
\[
  F(\{e_G\})Z=F(\{e_G\}X)=F(X)=Z,
\]
und ebenso \(ZF(\{e_G\})=Z\). Also ist \(F(\{e_G\})\)
neutral. Da ein neutrales Element eindeutig ist, folgt
\(F(\{e_G\})=\{e_H\}\).

Nun erhalten Isomorphismen auch Einheiten: Aus \(XY=YX=\{e_G\}\)
folgt durch Anwendung von \(F\)
\[
  F(X)F(Y)=F(Y)F(X)=\{e_H\}.
\]
Für die Umkehrabbildung gilt dasselbe. Mit
\eqref{reconstruction-reading-singletons} folgt deshalb, dass
\(F\) die Einermengen von \(G\) bijektiv auf die Einermengen
von \(H\) abbildet. Dies ist
\FormulaRefAuto{PowerGroupsSingletonLayer}[thm].

Für jedes \(g\in G\) gibt es somit genau ein \(\varphi(g)\in H\)
mit
\begin{equation}
  F(\{g\})=\{\varphi(g)\}.
  \tag{T}\label{reconstruction-reading-transport}
\end{equation}
Die drei Schritte lassen sich als Weg lesen:
\[
  g\ \longmapsto\ \{g\}
    \ \xrightarrow{\ F\ }\ \{\varphi(g)\}
    \ \longmapsto\ \varphi(g).
\]
Wir gehen vom Element zu seiner Einermenge, wenden den gegebenen
Isomorphismus an und nehmen anschließend das einzige Element
der Bildmenge. So entsteht eine Abbildung
\(\varphi\colon G\to H\).

Sie ist injektiv: Aus \(\varphi(g)=\varphi(h)\) folgt
\(F(\{g\})=F(\{h\})\), wegen der Injektivität von \(F\)
dann \(\{g\}=\{h\}\) und schließlich \(g=h\).
Sie ist surjektiv: Für \(k\in H\) ist \(\{k\}\) eine Einheit
von \(\PowNE H\). Ihr Urbild unter \(F\) ist daher eine
Einheit von \(\PowNE G\), also eine Einermenge \(\{g\}\).
Dann ist \(\varphi(g)=k\).

Es bleibt die Multiplikation. Für \(g,h\in G\) rechnen wir
\[
  \begin{aligned}
    \{\varphi(gh)\}
       &=F(\{gh\})\\
       &=F(\{g\}\{h\})\\
       &=F(\{g\})F(\{h\})\\
       &=\{\varphi(g)\}\{\varphi(h)\}
        =\{\varphi(g)\varphi(h)\}.
  \end{aligned}
\]
Zwei Einermengen sind genau dann gleich, wenn ihre Elemente
gleich sind. Also gilt \(\varphi(gh)=\varphi(g)\varphi(h)\).
Damit ist \(\varphi\) ein Gruppenisomorphismus. Er bildet auch
\(e_G\) auf \(e_H\) ab; dies folgt unmittelbar aus
\eqref{reconstruction-reading-transport} für das neutrale Element.
Der erste Rekonstruktionssatz ist vollständig bewiesen:
\FormulaRefAuto{PowerGroupsSingletonTransport}[thm]
(\ReconstructionProofLink{PowerGroupsSingletonTransport}).

\par\Needspace{13\baselineskip}
\section*{Wenn die zweite Struktur zunächst nur eine Halbgruppe ist}
\hypertarget{reconstruction.reading.identity}{}

Wir ändern jetzt ausdrücklich die Voraussetzung. Weiterhin ist
\(G\) eine Gruppe; die Zielstruktur \(S\) sei aber zunächst
nur eine Halbgruppe. Gegeben sei ein Isomorphismus
\[
  F\colon\PowNE G\longrightarrow\PowNE S.
\]
Hat \(S\) überhaupt ein neutrales Element? Und falls ja:
Weshalb sollte jedes seiner Elemente invertierbar sein?
Beides gehört nun zum Beweisziel.

Die erste Frage beantwortet bereits das Bild
\(E=F(\{e_G\})\). Als Element von \(\PowNE S\) ist \(E\)
eine nichtleere Teilmenge von \(S\); damit ist auch \(S\)
nichtleer. Die Neutralität von \(E\) zeigen wir ausdrücklich,
ohne bereits ein neutrales Element von \(S\) vorauszusetzen.
Sei \(Z\) irgendeine nichtleere Teilmenge von \(S\). Wegen der
Surjektivität gibt es \(X\in\PowNE G\) mit \(F(X)=Z\). Somit
\[
  \begin{aligned}
    EZ&=F(\{e_G\})F(X)=F(\{e_G\}X)=F(X)=Z,\\
    ZE&=F(X)F(\{e_G\})=F(X\{e_G\})=F(X)=Z.
  \end{aligned}
\]
Genau das besagt, dass \(E\) ein neutrales Element der
Potenzhalbgruppe ist. Noch wissen wir nicht, ob \(E\) eine
Einermenge ist. Für jedes \(s\in S\) dürfen wir aber
\(Z=\{s\}\) einsetzen und erhalten
\[
  E\{s\}=\{s\}=\{s\}E.
\]
Wählen wir ein \(u\in E\), so liegen \(us\) und \(su\)
jeweils in der Einermenge \(\{s\}\). Folglich ist
\(us=s=su\) für jedes \(s\in S\). Also ist \(u\) ein
neutrales Element der ursprünglichen Halbgruppe \(S\).

Dasselbe gilt für jedes weitere \(v\in E\). Dann folgt
\(uv=v\), weil \(u\) neutral ist, und \(uv=u\), weil \(v\)
neutral ist. Somit ist \(u=v\), also \(E=\{u\}\).
Wir schreiben fortan \(e_S=u\). Die Halbgruppe \(S\) ist
damit ein Monoid. Das ist der Gedanke hinter
\FormulaRefAuto{PowerSemigroupIdentityCharacterization}[thm]
aus Band~38.

Jetzt dürfen wir von den Einheiten von \(S\) sprechen.
Die bisherige Einermengenmethode erfasst zunächst nur diese
invertierbaren Elemente. Sie liefert noch keinen Grund dafür,
dass es in \(S\) keine anderen Elemente gibt. Um genau diesen
fehlenden Schritt zu begründen, untersuchen wir die ganze
Einheitenmenge als \emph{ein} Element der Potenzhalbgruppe.

\par\Needspace{13\baselineskip}
\section*{Die gesamte Einheitenmenge wiederfinden}
\hypertarget{reconstruction.reading.unitset}{}

Wir arbeiten für diesen Schritt mit zwei beliebigen Monoiden
\(A,B\) und einem Isomorphismus
\(F\colon\PowNE A\to\PowNE B\). Ihre Einheitenmengen heißen
\[
  \begin{aligned}
    U=A^\times&=\{u\in A\mid u\text{ ist invertierbar}\},\\
    V=B^\times&=\{v\in B\mid v\text{ ist invertierbar}\}.
  \end{aligned}
\]
Das hochgestellte Kreuz bezeichnet hier also die Einheitenmenge,
kein kartesisches Produkt. Beide Mengen sind nichtleer, weil
das jeweilige neutrale Element zu ihnen gehört.

Die Menge \(U\) bildet mit der aus \(A\) übernommenen
Multiplikation eine Gruppe, die \emph{Einheitengruppe} von
\(A\). Entsprechend ist \(V\) die Einheitengruppe von \(B\).
Wir wählen also die invertierbaren Elemente aus und behalten
ihre bisherige Multiplikation bei. Die
Assoziativität gilt bereits im Monoid, das neutrale Element
gehört dazu, und das Inverse einer Einheit ist wieder eine
Einheit. Auch das Produkt zweier Einheiten ist invertierbar:
Für \(u,v\in U\) hat \(uv\) das Inverse \(v^{-1}u^{-1}\),
denn
\[
  \begin{aligned}
    (uv)(v^{-1}u^{-1})&=u(vv^{-1})u^{-1}=e_A,\\
    (v^{-1}u^{-1})(uv)&=v^{-1}(u^{-1}u)v=e_A.
  \end{aligned}
\]
Der Satz, dass die Einheitenmenge eines Monoids mit dessen
Multiplikation eine Gruppe bildet, ist in Band~40 unter
\FormulaRefAuto{MonoidUnitSetIsGroup}[thm] bewiesen.

Unser Ziel lautet
\begin{equation}
  F(U)=V.
  \tag{U}\label{reconstruction-reading-unitset}
\end{equation}
Hier wird \(F\) einmal auf die \emph{Menge} \(U\) angewendet.
Das ist zulässig, weil \(U\) eine nichtleere Teilmenge von \(A\)
und daher ein einzelnes Element von \(\PowNE A\) ist.
Es geht an dieser Stelle weder um die Bilder aller Einermengen
noch um das Bild einer ganzen Familie von Teilmengen.

\medskip\noindent\textbf{Stabilität unter Einheiten.}
Wir nennen eine nichtleere Teilmenge \(X\subseteq A\)
\emph{einheitenstabil}, wenn für jede Einheit \(u\in U\)
\[
  \{u\}X=X=X\{u\}
\]
gilt. Multiplikation von links oder rechts mit einer Einheit
verändert also die Menge \(X\) als Ganzes nicht; ihre einzelnen
Elemente dürfen dabei vertauscht werden.

Zum Beispiel ist im Monoid der ganzen Zahlen mit der gewöhnlichen
Multiplikation \(U=\{-1,1\}\). Die Menge \(X=\{-2,2\}\)
ist einheitenstabil: Multiplikation mit \(1\) lässt beide
Elemente fest, Multiplikation mit \(-1\) vertauscht sie.
In beiden Fällen erhalten wir wieder genau \(X\).
Einheitenstabilität verlangt also nicht, dass die Elemente
von \(X\) selbst Einheiten sind oder dass \(X\) eine Gruppe bildet.

Da die Einheiten der Potenzhalbgruppe
genau die Einermengen der Einheiten von \(A\) sind, lässt sich
dies allein in ihrer Rechensprache ausdrücken: Jedes invertierbare
Element der Potenzhalbgruppe lässt \(X\) beim Multiplizieren
von beiden Seiten unverändert.

Deshalb wird Einheitenstabilität von \(F\) erhalten. Ist \(X\)
einheitenstabil und \(v\in V\), so ist \(\{v\}\) eine Einheit der
Zielpotenzhalbgruppe. Ihr Urbild ist eine Einheit, also eine
Einermenge \(\{u\}\) mit \(u\in U\). Daher gilt
\[
  \{v\}F(X)=F(\{u\})F(X)=F(\{u\}X)=F(X),
\]
und ebenso \(F(X)\{v\}=F(X)\). Dasselbe Argument gilt für
\(F^{-1}\). Dies erklärt
\FormulaRefAuto{SemigroupIsoUnitStabilityTransport}[thm].

Die Einheitenmenge \(U\) selbst ist einheitenstabil.
Für \(u\in U\) liegt jedes Produkt \(ux\) mit \(x\in U\)
wieder in \(U\). Umgekehrt lässt sich jedes \(w\in U\) als
\(w=u(u^{-1}w)\) schreiben, wobei \(u^{-1}w\in U\).
Also ist \(\{u\}U=U\). Entsprechend ergibt
\(w=(wu^{-1})u\) die Gleichheit \(U\{u\}=U\).
Dasselbe gilt für \(V\) im zweiten Monoid; siehe
\FormulaRefAuto{PowerMonoidUnitSetStable}[thm].

Einheitenstabilität allein zeichnet \(U\) jedoch noch nicht
eindeutig aus, wie das Beispiel \(\{-2,2\}\) zeigt.
Aus der Einheitenstabilität von \(F(U)\) allein
dürfen wir deshalb noch nicht \(F(U)=V\) folgern.

\medskip\noindent\textbf{Einheitenstabile Mengen und die ganze Einheitenmenge.}
Ist \(X\) einheitenstabil, dann gilt
\begin{equation}
  UX=X=XU.
  \tag{A}\label{reconstruction-reading-absorption}
\end{equation}
Denn jedes Element von \(UX\) hat die Form \(ux\) mit
\(u\in U\) und \(x\in X\). Wegen \(\{u\}X=X\) liegt
dieses Produkt in \(X\). Umgekehrt ist jedes \(x\in X\)
gleich \(e_Ax\) und liegt daher in \(UX\). Also ist
\(UX=X\). Die Gleichheit \(XU=X\) folgt auf dieselbe Weise
aus \(X\{u\}=X\) für alle \(u\in U\) und \(x=xe_A\).
Das ist \FormulaRefAuto{PowerMonoidUnitStableAbsorption}[thm].

Jetzt verwenden wir beide Richtungen des Isomorphismus.
Weil \(U\) einheitenstabil ist, ist \(F(U)\) einheitenstabil; im Zielmonoid
liefert \eqref{reconstruction-reading-absorption} daher
\[
  F(U)V=F(U).
\]
Weil \(V\) einheitenstabil ist, ist auch \(X=F^{-1}(V)\) einheitenstabil.
Im Ausgangsmonoid gilt somit \(UX=X\). Wenden wir darauf
\(F\) an, erhalten wir
\[
  F(U)V=F(U)F(X)=F(UX)=F(X)=V.
\]
Dasselbe Produkt \(F(U)V\) ist also einerseits \(F(U)\) und
andererseits \(V\). Das beweist
\eqref{reconstruction-reading-unitset}. Der Einsatz von
\(F^{-1}\) liefert genau die zusätzliche Information, die
uns bei der bloßen Einheitenstabilität fehlte. Die Aussage heißt in
Band~40 \FormulaRefAuto{LargePowerUnitSetImage}[thm]
(\ReconstructionProofLink{LargePowerUnitSetImage}).

\par\Needspace{13\baselineskip}
\section*{Von einer Menge zu allen ihren nichtleeren Teilmengen}
\hypertarget{reconstruction.reading.subsets}{}

Aus \(F(U)=V\) können wir nicht einfach schließen, dass \(F\)
jede Teilmenge von \(U\) in eine Teilmenge von \(V\) überführt.
Inklusionstreue gehört weiterhin nicht zu unseren Voraussetzungen.
Wir brauchen eine Bedingung, die das Enthaltensein durch ein
Mengenprodukt ausdrückt. Für jede nichtleere Teilmenge
\(Y\subseteq A\) gilt
\begin{equation}
  Y\subseteq U\quad\Longleftrightarrow\quad YU=U.
  \tag{K}\label{reconstruction-reading-criterion}
\end{equation}

Ist \(Y\subseteq U\), so liegen alle Produkte aus \(YU\)
wieder in der Gruppe \(U\). Also ist \(YU\subseteq U\).
Für die umgekehrte Inklusion wählen wir ein \(y_0\in Y\).
Jedes \(u\in U\) lässt sich schreiben als
\[
  u=y_0(y_0^{-1}u).
\]
Dabei ist \(y_0\in Y\) und \(y_0^{-1}u\in U\), also
\(u\in YU\). Damit ist \(YU=U\).

Gilt umgekehrt \(YU=U\), so liegt für jedes \(y\in Y\)
das Produkt \(ye_A=y\) in \(YU=U\). Folglich ist
\(Y\subseteq U\). Damit ist das Kriterium in beiden
Richtungen bewiesen; es ist
\FormulaRefAuto{PowerMonoidUnitSubsetCriterion}[thm].
Die erste Richtung enthält zugleich das allgemeine
Gruppenargument: Für jede nichtleere Teilmenge \(Y\) einer
Gruppe \(U\) ist \(YU=U=UY\). Die rechte Gleichheit folgt
entsprechend aus \(u=(uy_0^{-1})y_0\). Dies ist
\FormulaRefAuto{GroupPowerFullCarrierAbsorption}[thm].

Nun lässt sich das Produktkriterium durch \(F\) transportieren.
Für nichtleeres \(Y\subseteq A\) gilt die vollständige Kette
\[
  \begin{aligned}
    Y\subseteq U
      &\ \Longleftrightarrow\ YU=U\\
      &\ \Longleftrightarrow\ F(YU)=F(U)\\
      &\ \Longleftrightarrow\ F(Y)V=V\\
      &\ \Longleftrightarrow\ F(Y)\subseteq V.
  \end{aligned}
\]
Die zweite Äquivalenz benutzt die Injektivität von \(F\).
Die dritte benutzt die Multiplikationserhaltung und
\(F(U)=V\). Die letzte ist dasselbe Produktkriterium im
Monoid \(B\). In jeder Zeile sind die betrachteten Mengen
nichtleer, sodass alle Produkte in den Potenzhalbgruppen
definiert sind.

Damit erhalten wir nicht bloß eine Aussage über einzelne
Teilmengen, sondern über die ganze Familie:
\begin{equation}
  F[\PowNE U]=\PowNE V.
  \tag{F}\label{reconstruction-reading-family}
\end{equation}
Für die Inklusion von links nach rechts wenden wir die eben
bewiesene Kette auf jedes nichtleere \(Y\subseteq U\) an.
Für die andere Richtung sei \(Z\subseteq V\) nichtleer.
Wegen der Surjektivität von \(F\) hat \(Z\) ein Urbild
\(Y\in\PowNE A\). Aus \(F(Y)=Z\subseteq V\) folgt mit
der Kette rückwärts \(Y\subseteq U\). Somit liegt \(Z\)
tatsächlich im Familienbild auf der linken Seite.

Die Schreibweisen sollten wir sorgfältig auseinanderhalten:
\(F(U)\) ist das Bild einer einzelnen Menge, während
\[
  F[\PowNE U]=\{F(Y)\mid Y\in\PowNE U\}
\]
eine Familie von Mengen ist. Das erste Objekt ist eine
Teilmenge von \(B\), das zweite eine Teilmenge von
\(\PowNE B\). Die Gleichheit der ersten
Objekte wurde mit Einheitenstabilität bewiesen, die Gleichheit der
Familien anschließend mit dem Produktkriterium. Gemeinsam
ergeben sie
\FormulaRefAuto{LargePowerSemigroupUnitGroupReconstruction}[thm].

\par\Needspace{13\baselineskip}
\section*{Warum eine einzige bekannte Gruppe genügt}
\hypertarget{reconstruction.reading.rigidity}{}

Wir kehren zu unserem stärkeren Ziel zurück. Es seien \(G\)
eine Gruppe, \(S\) eine Halbgruppe und
\(F\colon\PowNE G\to\PowNE S\) ein Isomorphismus.
Schon bewiesen ist, dass \(S\) ein neutrales Element besitzt,
also ein Monoid ist. Wir dürfen daher die Ergebnisse der
letzten beiden Abschnitte anwenden.

Die Einheitenmenge von \(G\) ist ganz \(G\). Nennen wir die
Einheitenmenge von \(S\) wieder \(V=S^\times\), so liefert
\eqref{reconstruction-reading-family}
\[
  F[\PowNE G]=\PowNE V.
\]
Andererseits ist \(F\) auf ganz \(\PowNE S\) surjektiv.
Deshalb gilt zugleich \(F[\PowNE G]=\PowNE S\), also
\[
  \PowNE S=\PowNE V.
\]
Für jedes \(s\in S\) liegt somit \(\{s\}\) in \(\PowNE V\).
Das bedeutet \(s\in V\). Da ohnehin \(V\subseteq S\) gilt,
folgt \(S=V\): Jedes Element von \(S\) ist invertierbar.
Die Zielhalbgruppe ist also eine Gruppe.

Jetzt sind wir in der Situation des bereits bewiesenen ersten
Rekonstruktionssatzes. Es gibt einen Gruppenisomorphismus
\(\varphi\colon G\to S\), eindeutig festgelegt durch
\[
  F(\{g\})=\{\varphi(g)\}\qquad(g\in G).
\]
Wir haben damit den stärkeren Satz
\FormulaRefAuto{GroupPowerSemigroupRigidity}[thm] erreicht:
\begin{quote}
  Ist die große Potenzhalbgruppe einer Halbgruppe zu derjenigen
  einer Gruppe isomorph, so ist die Halbgruppe selbst eine
  Gruppe und zu dieser Gruppe isomorph.
\end{quote}
Die formale Schlussableitung ist ebenfalls ausgelagert
(\ReconstructionProofLink{GroupPowerSemigroupRigidity}).
Der zusätzliche Gedanke gegenüber dem ersten Satz war nicht
eine andere Konstruktion von \(\varphi\). Er bestand darin,
über die Einheitenmenge und ihre Teilmengenfamilie sämtliche
Elemente der Zielstruktur zu erfassen.

\par\Needspace{13\baselineskip}
\section*{Was das Ergebnis sagt und wo der Beweis seine Grenze hat}
\hypertarget{reconstruction.reading.limits}{}

In keinem Schritt wurden Endlichkeit oder Kommutativität
vorausgesetzt. Auch die Gruppe mit nur einem Element ist
eingeschlossen. Die Reihenfolge der Faktoren bleibt in allen
Rechnungen erhalten; insbesondere steht beim Inversen von
\(uv\) die umgekehrte Reihenfolge \(v^{-1}u^{-1}\).

Dass wir die \emph{große} Potenzhalbgruppe betrachten, ist
für die zweite Beweisstufe wesentlich. Die gesamte
Einheitenmenge \(U\) muss selbst eine zulässige Eingabe von
\(F\) sein. Beschränkte man sich auf endliche nichtleere
Teilmengen, könnte bei einer unendlichen Einheitengruppe
\(F(U)\) gar nicht gebildet werden. Der hier geführte
stärkere Beweis lässt sich deshalb nicht unverändert auf
diese kleinere Potenzhalbgruppe übertragen. Damit ist keine
gegenteilige Behauptung über sie bewiesen; es ist die genaue
Grenze unseres Arguments.

Die erste Beweisstufe, bei der beide Ausgangsstrukturen
bereits Gruppen sind, kommt hingegen nur mit Einermengen,
ihren Inversen und Produkten aus. Sie funktioniert auch
für die Potenzhalbgruppen der endlichen nichtleeren
Teilmengen. Die beiden Beweisstufen unterscheiden sich
also nicht nur im Ziel, sondern auch in der Information,
die sie verwenden.

Ebenso wenig haben wir behauptet, dass der gegebene
Isomorphismus \(F\) auf jeder Teilmenge punktweise durch
\(\varphi\) entsteht. Bewiesen ist
\(F(\{g\})=\{\varphi(g)\}\). Die weitergehende Gleichung
\[
  F(X)=\varphi[X]
       =\{\varphi(x)\mid x\in X\}
\]
für beliebige nichtleere \(X\subseteq G\) folgt nicht aus
unserer Konstruktion und wird für die Rekonstruktion der
Gruppe nicht benötigt. Insbesondere haben wir keine
allgemeine Erhaltung von Inklusionen benutzt. Für Teilmengen
der Einheitenmenge wurde das benötigte Enthaltensein eigens
durch die Gleichung \(YU=U\) erkannt.

Hier liegt auch der Unterschied zum Mogiljanskaja-Gegenbeispiel
aus Band~28. Bei allgemeinen Halbgruppen können die großen
Potenzhalbgruppen isomorph sein, obwohl die Ausgangsstrukturen
es nicht sind. Die Einermengen müssen dann nicht als solche
wiedererkennbar bleiben. Bei Gruppen zeichnet die
Invertierbarkeit gerade diese Einermengen aus; die zusätzliche
Betrachtung der Einheitenmenge reicht sogar aus, um eine
zunächst unbekannte Halbgruppe als Gruppe zu erkennen.

\clearpage
\section*{Welche Bijektion verlangt die Potenzhalbgruppenvermutung?}
\hypertarget{reconstruction.reading.conjecture}{}

Nun lässt sich unsere Rekonstruktion in die weitergehende
Frage aus Band~37 einordnen. Die dort formulierte
\emph{endliche Potenzhalbgruppenvermutung} besagt: Für
endliche Halbgruppen \(S,T\) soll gelten
\[
  \PowNE S\cong\PowNE T\quad\Longrightarrow\quad S\cong T.
\]
Das Zeichen \(\cong\) bedeutet jeweils, dass ein Isomorphismus
existiert. Aus irgendeinem Potenzisomorphismus \(F\) soll also
die Existenz eines Grundisomorphismus \(f\) folgen.
Die allgemeine endliche Aussage ist eine offene Vermutung;
der hier bewiesene Gruppenfall ist ein Spezialfall davon.\footnote{Zur
offenen endlichen Frage siehe Liu--Tringali,
\href{https://arxiv.org/html/2606.01917v1}{\emph{Power Semigroups and Two
Rigidity Theorems for Groups}} (2026), Einleitung, Questions~1.1.
Ohne Endlichkeit ist die allgemeine Behauptung durch das
Mogiljanskaja-Gegenbeispiel widerlegt.}

\medskip\noindent\textbf{Eine andere Bijektion ist zulässig,
aber sie muss die Produkte erhalten.}
Gesucht ist eine bijektive Abbildung \(f\colon S\to T\), für die
\[
  f(st)=f(s)f(t)\qquad(s,t\in S)
\]
gilt. Wie man diese Abbildung findet, schreibt die Vermutung
nicht vor. Eine beliebige Bijektion reicht allerdings nicht.
Im endlichen Fall liefert schon
\(2^{|S|}-1=2^{|T|}-1\) die Gleichheit \(|S|=|T|\), also
irgendeine Bijektion. Offen bleibt gerade, ob es eine gibt,
die auch die Multiplikation richtig zuordnet.

Unser Einermengenweg ist eine hinreichende Methode: Bildet
\(F\) die Einermengen von \(S\) bijektiv auf die von \(T\)
ab, definiert \(F(\{s\})=\{f(s)\}\) den gewünschten
Isomorphismus. Dafür genügt dieselbe Rechnung mit
\(\{s\}\{t\}=\{st\}\) wie im Gruppenbeweis.
Die Vermutung verlangt aber \emph{nicht}, dass jeder
vorgegebene Potenzisomorphismus diese Methode ermöglicht.

\medskip\noindent\textbf{Ein kleines Beispiel zeigt den Unterschied.}
Sei \(S=\{0,a\}\) mit \(xy=0\) für alle \(x,y\in S\).
Das ist eine Halbgruppe: Beide Klammerungen eines Dreifachprodukts
ergeben \(0\). Jedes Produkt nichtleerer Teilmengen ist
\(\{0\}\). Definieren wir
\[
  F(\{0\})=\{0\},\qquad
  F(\{a\})=S,\qquad F(S)=\{a\},
\]
so vertauscht \(F\) zwei der drei Eingaben und ist bijektiv.
Außerdem gilt für alle nichtleeren \(X,Y\subseteq S\)
\[
  F(XY)=F(\{0\})=\{0\}=F(X)F(Y).
\]
Damit ist \(F\) ein Potenzhalbgruppenisomorphismus, der die
Einermenge \(\{a\}\) auf eine Zweiermenge schickt. Aus
diesem Bild lässt sich kein einziges \(f(a)\) ablesen.
Trotzdem ist \(f=\mathrm{id}_S\) ein Isomorphismus der
Grundhalbgruppe. Sein induzierter Potenzisomorphismus
\(X\mapsto f[X]=X\) ist ein anderer als \(F\).
Das Beispiel widerlegt die Forderung, dass \emph{jedes}
\(F\) Einermengen erhalten müsse; es widerlegt nicht die
Potenzhalbgruppenvermutung.

\medskip\noindent\textbf{Bei Gruppen ist die Einermengenerhaltung
für jedes gegebene \(F\) erzwungen.}
Hier sind die Einermengen genau die Einheiten. Deshalb kann
kein Potenzisomorphismus diese Schicht verfehlen. Das zu
\(F\) gehörende \(\varphi\) ist durch die Bilder der
Einermengen eindeutig bestimmt. Man darf dennoch einen
anderen Beweisweg wählen oder, sofern vorhanden, einen
anderen Gruppenisomorphismus verwenden. Dieser muss nicht
dieselben Einermengenbilder wie das vorgegebene \(F\) liefern.
Erhalten bleibt die Unterscheidung aus dem vorigen Abschnitt:
Auch hier folgt aus der Konstruktion keine Gleichung
\(F(X)=\varphi[X]\) für alle nichtleeren Teilmengen \(X\).

Falls wir umgekehrt einen Potenzisomorphismus suchen, um
ein Gegenbeispiel zur Vermutung zu bauen, genügt daher
kein bloßes Vertauschen von Einermengen und größeren Mengen.
Zusätzlich müssten die Grundhalbgruppen \(S,T\) nachweislich
\emph{nicht} isomorph sein. Mit einer Gruppe auf einer Seite
ist ein solches Gegenbeispiel für die große Potenzhalbgruppe
durch unseren Satz ausgeschlossen. Bei allgemeinen
Halbgruppen kann man andere Bijektionen zwischen den
Teilmengen untersuchen; ihre Multiplikationserhaltung
bleibt stets zu beweisen.

Die zugehörigen Definitionen und Satzaussagen bleiben unter
\ReconstructionMainLink{} zugänglich. Die zwölf formalen
Ableitungstabellen dieses Lesebogens sind in den
\ReconstructionProofsLink{} zusammengeführt. Die allgemeinen
Grundlagen über Monoide und Einheitengruppen stehen weiterhin
in den Fachbänden; alle für das Verständnis tragenden
Argumente wurden hier im Text ausgeführt.

\clearpage
\section*{Historischer Abriss}
\hypertarget{reconstruction.reading.history}{}

Die Frage, wie viel eine Potenzhalbgruppe über ihre
Ausgangsstruktur verrät, wurde in den 1960er Jahren
systematisch untersucht. T.~Tamura und J.~Shafer
veröffentlichten 1967 ihre Arbeit \emph{Power semigroups}.
Ein damaliges Forschungsabstract nennt ausdrücklich den
Rekonstruktionssatz für zwei \emph{endliche} Gruppen.
Shafer behandelte noch im selben Jahr in einer kurzen
Notiz den Fall beliebiger Gruppen.\footnote{T.~Tamura und
J.~Shafer, \emph{Power semigroups}, \emph{Math. Japon.}~12
(1967), S.~25--32; J.~Shafer, \emph{Note on power semigroups},
ebenda, S.~32. Der endliche Satz steht auch im
\href{https://www.ams.org/journals/notices/196708/196708FullIssue.pdf}{Forschungsabstract 648-194},
\emph{Notices of the American Mathematical Society}~14 (1967),
S.~688. Zur Zuordnung des allgemeinen Zweigruppensatzes zu
Shafers Notiz siehe Liu--Tringali, Einleitung und Beweis von
Theorem~1.3, in der unten genannten Arbeit.}
Das ist die erste Beweisstufe unseres Textes: Beide
Ausgangsstrukturen sind schon als Gruppen bekannt, und
ihre Einermengen lassen sich an der Invertierbarkeit
wiedererkennen.

Für allgemeine Halbgruppen fällt die Antwort anders aus.
E.~M.~Mogiljanskaja veröffentlichte 1973 ein Gegenbeispiel:
Zwei nichtisomorphe Halbgruppen können isomorphe
Potenzhalbgruppen besitzen.\footnote{E.~M.~Mogiljanskaja,
\href{https://doi.org/10.1007/BF02389140}{\emph{Non-isomorphic semigroups with isomorphic semigroups of subsets}},
\emph{Semigroup Forum}~6 (1973), S.~330--333.}
Die Lesefassung zu Band~28 erläutert diese Grenze anhand
einer ausgeführten Konstruktion. Für die gegenwärtige Frage
ist die Folgerung entscheidend: Die bloße Existenz aller
Mengenprodukte garantiert noch keine Rekonstruktion.
Zusätzliche Eigenschaften der Ausgangsstruktur müssen
erklären, weshalb die ursprünglichen Elemente erkennbar
bleiben.

Der stärkere Satz dieser Lesefassung erscheint 2026 im
Preprint \emph{Power Semigroups and Two Rigidity Theorems
for Groups} von Shuolin Liu und Salvatore Tringali.
Ihr Theorem~1.3 setzt nur auf einer Seite eine Gruppe
voraus; auf der anderen Seite genügt eine Halbgruppe.
Ein zentrales Hilfsergebnis ist Theorem~2.3: Ein Isomorphismus
der großen Potenzhalbgruppen zweier Monoide transportiert
sowohl die gesamte Einheitenmenge als auch die Familie
ihrer nichtleeren Teilmengen.\footnote{Shuolin Liu und
Salvatore Tringali,
\href{https://arxiv.org/html/2606.01917v1}{\emph{Power Semigroups and Two Rigidity Theorems for Groups}},
arXiv:2606.01917v1, 1.~Juni 2026, insbesondere Theoreme~1.3
und~2.3 sowie der Beweis von Theorem~1.3 in Abschnitt~2.}
Genau diese beiden Aussagen haben wir als
\(F(U)=V\) und \(F[\PowNE U]=\PowNE V\) kennengelernt.

Die historische Unterscheidung entspricht damit der
logischen Gliederung des Beweises. Der ältere Satz erkennt
unter bereits bekannten Gruppen die richtige wieder.
Der stärkere Satz erkennt zusätzlich, dass die unbekannte
Zielhalbgruppe überhaupt eine Gruppe sein muss. Dazu
genügt es nicht, nur die invertierbaren Einermengen zu
transportieren; die gesamte Einheitenmenge vermittelt
den entscheidenden weiteren Schluss.

Unsere Darstellung ordnet diese Gedanken für den Einstieg
neu: Die Zweiergruppe liefert ein durchgehendes Bild,
die allgemeinen Monoidargumente werden einzeln begründet,
und erst danach wird der stärkere Satz zusammengesetzt.
Die historische Zuordnung der Resultate und diese
didaktische Ausarbeitung sind dabei zu unterscheiden.
Der Text gibt einen eigenständigen Lesebeweis der genannten
Aussagen, keine Abschrift der historischen Originalbeweise.

\par\Needspace{16\baselineskip}
\section*{Schlussbemerkung}
\hypertarget{reconstruction.reading.conclusion}{}

Am Anfang stand die kleine Multiplikationstafel der drei
nichtleeren Teilmengen einer Zweiergruppe. In ihr konnten
wir die beiden ursprünglichen Elemente wiederfinden, weil
ihre Einermengen genau die invertierbaren Eingaben waren.
Der allgemeine Beweis hat diesen Gedanken unverändert
beibehalten: Invertierbarkeit macht die Einermengen einer
Gruppe allein an den Rechengesetzen erkennbar.

Der stärkere Satz verlangte eine weitere Unterscheidung.
Die Einheiten als einzelne Elemente zu erkennen ist etwas
anderes, als ihre gesamte Menge wiederzufinden. Für diese
Menge verbanden wir die Stabilität unter Einheiten mit
dem Isomorphismus und seiner Umkehrabbildung. Das
Produktkriterium \(YU=U\) erschloss anschließend alle ihre
nichtleeren Teilmengen. Dadurch blieb in der Zielhalbgruppe
kein Element außerhalb der Einheiten übrig.

So enthält die Multiplikation aller nichtleeren Teilmengen
genug Information, um eine Gruppe bis auf Isomorphie zu
bestimmen, selbst wenn auf der anderen Seite zunächst nur
eine Halbgruppe vorausgesetzt wird. Der rekonstruierte
Gruppenisomorphismus ist an den Bildern der Einermengen
ablesbar. Dafür mussten wir weder eine Ordnung der
Teilmengen noch ihre Größen kennen: Die benötigte
Information lag bereits in ihren Produkten.

Für die allgemeine Potenzhalbgruppenvermutung ist die
Erkennung aller Einermengen damit ein möglicher Beweisweg,
keine zusätzliche Forderung an jeden vorgegebenen
Isomorphismus. Entscheidend ist am Ende eine Bijektion
der Grundträger, die ihre Multiplikation erhält.
